\chapter{İstatistiksel Düşünme ve İstatistiksel Araştırma Süreci}
\label{ch:research}

\begin{prismbox}[PrismLight]{Öğrenme çıktıları}
Bu bölümün sonunda okuyucu bir araştırma sorusunu değişkenlere dönüştürebilecek;
evren, örneklem, parametre ve istatistik kavramlarını ayırt edebilecek; veri
türüne uygun bir istatistiksel çalışma planı kurabilecektir.
\end{prismbox}

\begin{prerequisitebox}
Temel aritmetik işlemler, yüzde kavramı ve basit tablo okuma becerisi
yeterlidir. Bu bölüm sonraki tüm bölümlerin araştırma dili ve karar
çerçevesini kurar.
\end{prerequisitebox}

\section{İstatistik neden yalnızca hesaplama değildir?}

İstatistik, eldeki sayılara bir formül uygulamaktan önce verinin nasıl
üretildiğini sorgulayan bir düşünme biçimidir. Sağlam bir analiz şu zinciri
izler:
\[
\begin{aligned}
\text{araştırma sorusu}&\longrightarrow\text{ölçülebilir değişkenler}
\longrightarrow\text{veri toplama tasarımı}\\
&\longrightarrow\text{analiz}\longrightarrow\text{bağlamsal yorum}.
\end{aligned}
\]
Zincirin ilk halkalarında yapılan bir hata, daha sonra kullanılan yöntem ne
kadar gelişmiş olursa olsun giderilemez. Örneğin yalnız gönüllü öğrencilerden
toplanan bir memnuniyet anketinde örnekleme yanlılığı varsa, çok küçük bir
\pval-değeri bu yanlılığı ortadan kaldırmaz. Bu nedenle çağdaş istatistik
eğitimi hesaplamadan önce araştırma sorusuna ve veri üretim sürecine ağırlık
verir \citep{gaise2016}.

\section{Temel kavramlar}

\begin{definitionbox}
\begin{description}[style=nextline,leftmargin=2.8cm,labelwidth=2.5cm]
  \item[Evren (anakütle)] Araştırma sonucunun genellenmek istendiği bütün birimlerdir.
  \item[Örneklem] Evrenden gözlenen veya ölçülen birimler kümesidir.
  \item[Parametre] Evrene ait, çoğu zaman bilinmeyen sayısal özelliktir.
  \item[İstatistik] Örneklem verisinden hesaplanan sayısal özettir.
  \item[Değişken] Birimden birime farklı değer alabilen ölçülen özelliktir.
\end{description}
\end{definitionbox}

Bir üniversitedeki bütün birinci sınıf öğrencilerinin matematik kaygısı
ortalaması bir parametre, seçilen 120 öğrencinin ortalaması ise bir
istatistiktir. İstatistik, parametre hakkında bilgi edinmek için kullanılır.

\section{Değişken türleri ve ölçme düzeyleri}
\label{sec:research-variable-types}

\begin{table}[H]
\centering
\caption{Ölçme düzeyleri ve uygun temel özetler.}
\label{tab:measurement-levels}
\begin{tabular}{p{0.18\textwidth}p{0.31\textwidth}p{0.35\textwidth}}
\toprule
Düzey & Örnek & Uygun özetler \\
\midrule
Nominal & Program türü, cinsiyet kategorisi & Frekans, yüzde, mod \\
Sıralı & Memnuniyet düzeyi, Likert maddesi & Frekans, medyan, yüzdelikler \\
Aralık & Standartlaştırılmış test puanı & Ortalama, standart sapma, fark \\
Oran & Süre, yaş, doğru cevap sayısı & Tüm aritmetik özetler ve oranlar \\
\bottomrule
\end{tabular}
\end{table}

Sayısal kod verilmesi bir değişkeni otomatik olarak nicel yapmaz. ``Hiç
katılmıyorum'' ile ``tamamen katılıyorum'' seçeneklerine 1--5 kodlarının
verilmesi kategoriler arasında bir sıra oluşturur; fakat tek bir madde için
ardışık kategoriler arasındaki uzaklıkların kesinlikle eşit olduğu sonucunu
garanti etmez.

\begin{mistakebox}
Öğrenci numarası rakamlardan oluşsa da nicel bir değişken değildir. Bu
numaraların ortalamasını hesaplamak araştırma açısından anlam taşımaz.
\end{mistakebox}

\section{İstatistiksel araştırma döngüsü}
\label{sec:research-cycle}

\begin{enumerate}
  \item Araştırma problemi ve hedef evren açıkça tanımlanır.
  \item Değişkenler işlemsel olarak tanımlanır ve ölçme düzeyleri belirlenir.
  \item Örnekleme ve veri toplama yöntemi seçilir.
  \item Veri temizlenir; eksik, olanaksız ve aykırı değerler incelenir.
  \item Önce grafiksel ve betimsel, sonra uygun çıkarımsal analiz yapılır.
  \item Belirsizlik, etki büyüklüğü, sınırlılıklar ve bağlam birlikte raporlanır.
\end{enumerate}
Eksik verinin oluşma mekanizması ve kullanılan yöntemin açıkça raporlanması
için bkz. \citet{little2019}.

Şekil~\ref{fig:research-cycle} bu adımların birbirine nasıl bağlandığını
gösterir. Sonuçların yeni sorular doğurması döngüyü yeniden başlatır;
bu, aynı veride istenen sonucu bulana kadar analiz değiştirmek değildir.

\begin{figure}[H]
\centering
\begin{tikzpicture}[
  stage/.style={draw=PrismBlue,fill=PrismLight,rounded corners=2mm,
    text width=3.65cm,minimum height=1.35cm,align=center,
    font=\small,inner sep=5pt},
  flow/.style={->,draw=PrismBlue,line width=0.8pt}]
  \node[stage] (question) at (-2.55,0)
    {\textbf{1. Araştırma sorusu}\\Kim, ne, hangi evren?};
  \node[stage] (variables) at (2.55,0)
    {\textbf{2. Ölçülebilir değişkenler}\\Ne, nasıl ölçülecek?};
  \node[stage] (design) at (2.55,-1.8)
    {\textbf{3. Veri toplama tasarımı}\\Kimlerden, hangi yöntemle?};
  \node[stage] (cleaning) at (-2.55,-1.8)
    {\textbf{4. Veri kontrolü}\\Eksik, hatalı,\\aykırı değerler};
  \node[stage] (analysis) at (-2.55,-3.6)
    {\textbf{5. Analiz}\\Grafikler, özetler\\ve çıkarım};
  \node[stage] (report) at (2.55,-3.6)
    {\textbf{6. Bağlamsal yorum}\\Belirsizlik ve sınırlılıklar};
  \draw[flow] (question)--(variables);
  \draw[flow] (variables)--(design);
  \draw[flow] (design)--(cleaning);
  \draw[flow] (cleaning)--(analysis);
  \draw[flow] (analysis)--(report);
  \draw[flow,dashed] (report.east)--(5,-3.6)--(5,1.05)
    --(-2.55,1.05)--(question.north);
  \node[font=\footnotesize,fill=white,inner sep=3pt] at (0,1.05)
    {Yeni araştırma soruları};
\end{tikzpicture}
\caption{Araştırma sorusundan yoruma uzanan istatistiksel araştırma döngüsü.
Kesikli ok, bulgulardan yeni bir araştırmaya geçişi gösterir.}
\label{fig:research-cycle}
\end{figure}

\section{Örnekleme ve yanlılık}
\label{sec:research-bias}

Basit rastgele örneklemede her birimin bilinen bir seçilme şansı vardır.
Tabakalı örneklemede önemli alt gruplar ayrı ayrı temsil edilir. Küme
örneklemesinde doğal gruplar, örneğin sınıflar veya okullar seçilir. Kolayda
örnekleme hızlıdır; ancak hedef evreni temsil etme gücü çoğu zaman sınırlıdır.

\begin{researchbox}
Yanlılık örneklem hacmi büyütülerek otomatik olarak giderilemez. Yanlı bir
seçim mekanizmasından gelen 10.000 gözlem, iyi tasarlanmış daha küçük bir
örneklemden daha yanıltıcı olabilir.
\end{researchbox}

\section{Etik ve eleştirel veri okuryazarlığı}

İstatistiksel sonuçlar kişisel verilerin korunması, katılımcı onamı, eksik
verilerin açıklanması ve analiz kararlarının şeffaflığıyla birlikte
değerlendirilir. Yalnız anlamlı sonuçları seçerek raporlamak, korelasyonu
nedensellik gibi sunmak veya analiz sonrasında hipotezi önceden belirlenmiş
gibi göstermek bilimsel güvenilirliği zedeler
\citep{asaethics2022,appelbaum2018}.

\section{Çözümlü örnek: araştırma sorusundan plana}

``Derse devam süresi başarı puanıyla ilişkili midir?'' sorusunda gözlem birimi
öğrenci, açıklayıcı değişken haftalık devam süresi, yanıt değişkeni başarı
puanıdır. Hedef evren fakültedeki ilgili dersi alan tüm öğrencilerdir. Önce
iki nicel değişken için saçılım grafiği ve korelasyon düşünülebilir. Ancak
çalışma gözlemselse bulunan ilişki, devam süresinin başarıya neden olduğunu
tek başına göstermez; önceki başarı ve çalışma alışkanlığı gibi değişkenler
iki sonucu birlikte etkileyebilir \citep{shadish2002}.

\section{Yazılım uygulaması: veri yapısını tanıma}

\begin{softwarebox}
\begin{lstlisting}
import pandas as pd

data = pd.DataFrame({
    "devam_saati": [8, 10, 7, 12, 9],
    "basari": [68, 75, 64, 83, 72],
    "program": ["A", "A", "B", "B", "A"]
})

print(data.info())
print(data.describe(include="all"))
\end{lstlisting}
İlk çıktı değişken türlerini ve eksik değerleri, ikinci çıktı temel özetleri
gösterir. Yazılımın atadığı türler araştırmadaki ölçme düzeyiyle ayrıca
karşılaştırılmalıdır.
\end{softwarebox}

\begin{assumptionbox}
İstatistiksel analizden önce her satırın hangi gözlem birimini, her sütunun
hangi değişkeni temsil ettiği ve aynı birimin yinelenip yinelenmediği
doğrulanmalıdır. Belirsiz veri yapısı doğru yöntem seçimini engeller.
\end{assumptionbox}

\section{Çözümlü problemler}

Bu problemler araştırma sorusunu hedef evren, gözlem birimi, değişken,
örnekleme tasarımı ve uygun ilk analiz kararlarına dönüştürme becerisini
geliştirir. Her çözümde yalnız sonuç değil, kararın gerekçesi de gösterilir.

\subsection*{Problem 1: Araştırma sorusunun bileşenleri}

Bir araştırmacı, ``Ortaokul öğrencilerinin günlük ekran süresi ile matematik
başarısı arasında ilişki var mıdır?'' sorusunu incelemek istemektedir. Hedef
evreni, gözlem birimini, değişkenleri ve araştırma türünü belirleyin.

\textbf{Çözüm.} Sonucun genellenmek istendiği topluluk araştırmanın yapıldığı
bölgedeki ortaokul öğrencileridir; dolayısıyla hedef evren bu öğrencilerin
tamamıdır. Her bir öğrenci bir gözlem birimidir. Günlük ekran süresi saat veya
dakika olarak ölçülen nicel açıklayıcı değişken, matematik başarı puanı ise
nicel yanıt değişkenidir. Öğrencilerin mevcut alışkanlıkları değiştirilmeden
yalnız ölçüm yapılıyorsa çalışma gözlemseldir.

İlk incelemede iki değişken için saçılım grafiği çizilebilir ve ilişkinin
doğrusal olması durumunda korelasyon hesaplanabilir. Bununla birlikte gözlenen
ilişki nedensellik kanıtı değildir. Önceki başarı, aile desteği, uyku süresi
ve çalışma alışkanlığı gibi değişkenler hem ekran süresiyle hem başarıyla
ilişkili olabilir.

\textbf{Raporlama örneği.} ``Çalışma, bölgedeki ortaokul öğrencilerinde günlük
ekran süresi ile matematik başarı puanı arasındaki ilişkiyi gözlemsel bir
tasarımla incelemektedir. Bulgular ilişki göstergesi olarak yorumlanacak,
nedensel etki iddiasında bulunulmayacaktır.''

\subsection*{Problem 2: Ölçme düzeyi ve uygun özet}

Bir öğrenci anketinde program türü, sınıf düzeyi, bir maddelik memnuniyet
yanıtı, haftalık çalışma süresi ve sınav puanı kaydedilmiştir. Değişkenlerin
ölçme düzeylerini ve uygun ilk özetlerini belirleyin.

\textbf{Çözüm.}

\begin{center}
\small
\begin{tabularx}{\textwidth}{@{}p{0.25\textwidth}p{0.22\textwidth}X@{}}
\toprule
Değişken & Ölçme düzeyi & Uygun ilk özet \\
\midrule
Program türü & Nominal & Frekans ve yüzde \\
Sınıf düzeyi & Sıralı & Frekans, yüzde ve medyan \\
Memnuniyet maddesi & Sıralı & Kategori yüzdeleri ve medyan \\
Haftalık çalışma süresi & Oran & Ortalama, medyan, standart sapma ve histogram \\
Sınav puanı & Aralık/nicel & Ortalama, standart sapma ve dağılım grafiği \\
\bottomrule
\end{tabularx}
\end{center}

Program türüne 1, 2 ve 3 kodları verilmiş olsa bile bu kodların ortalaması
anlamlı değildir. Benzer biçimde tek bir Likert maddesinde 1 ile 2 arasındaki
farkın 4 ile 5 arasındaki farka kesin olarak eşit olduğu varsayılmamalıdır.
Çalışma süresinde gerçek bir sıfır noktası bulunduğundan ``sekiz saat, dört
saatin iki katıdır'' yorumu anlamlıdır.

\textbf{Sık hata.} Yazılımın bir sütunu tam sayı olarak tanıması, o değişkenin
araştırma bakımından nicel olduğunu göstermez. Ölçme düzeyi kod biçiminden
değil, değerlerin taşıdığı anlamdan belirlenir.

\subsection*{Problem 3: Parametre ile istatistiği ayırma}

Bir üniversitede 6.000 lisans öğrencisi bulunmaktadır. Rastgele seçilen 150
öğrencinin haftalık çalışma süresi ortalaması 7,2 saat, standart sapması 2,4
saat bulunmuştur. Bu araştırmada evreni, örneklemi, parametreyi ve
istatistikleri belirleyin.

\textbf{Çözüm.} Evren 6.000 lisans öğrencisinin tamamı, örneklem ise ölçüm
yapılan 150 öğrencidir. Bütün öğrencilerin bilinmeyen ortalama çalışma süresi
$\mu$ bir parametredir. Örneklem ortalaması $\bar{x}=7{,}2$ ve örneklem
standart sapması $s=2{,}4$ ise örnekleme bağlı olarak değişebilen
istatistiklerdir.

7,2 saatlik değer doğrudan ``üniversitedeki öğrencilerin gerçek ortalaması''
diye sunulmamalıdır. Bu değer $\mu$ için bir nokta tahminidir ve örnekleme
değişkenliği içerir. Ayrıca rastgele seçildiği belirtilen öğrencilerin hangi
listeden ve nasıl seçildiği de temsil gücü açısından açıklanmalıdır.

\textbf{Raporlama örneği.} ``Rastgele seçilen 150 öğrencinin haftalık çalışma
süresi ortalaması 7,2 saat ($s=2{,}4$) bulunmuştur. Örneklem ortalaması,
üniversitedeki 6.000 lisans öğrencisinin bilinmeyen ortalamasını tahmin etmek
için kullanılacaktır.''

\subsection*{Problem 4: Örnekleme tasarımı ve yanlılık}

Bir üniversite yönetimi uzaktan eğitim memnuniyetini ölçmek için anket
bağlantısını yalnız üniversitenin sosyal medya hesabında paylaşmış ve üç gün
içinde gönüllü olarak yanıt veren 2.400 öğrencinin sonuçlarını incelemiştir.
Örneklem büyüktür. Buna rağmen hangi sorunlar ortaya çıkabilir ve tasarım
nasıl iyileştirilebilir?

\textbf{Çözüm.} Anketi yalnız sosyal medya hesabını izleyen öğrenciler
görebildiği için örnekleme çerçevesi hedef evreni tam kapsamayabilir. Bu durum
\emph{kapsam yanlılığı} doğurur. Katılım gönüllü olduğu için çok memnun veya
çok memnuniyetsiz öğrenciler yanıt vermeye daha istekli olabilir; bu da
\emph{gönüllülük yanlılığı} oluşturur. Ayrıca üç günlük süre, hesabı seyrek
kontrol eden öğrencileri dışlayabilir.

Öğrenci bilgi sistemindeki güncel kayıt listesinden olasılıklı bir örneklem
seçmek daha uygun olur. Fakülte veya sınıf düzeylerinin yeterince temsil
edilmesi önemliyse tabakalı örnekleme kullanılabilir. Seçilen öğrencilere
birden fazla hatırlatma gönderilmeli, yanıt oranı ve yanıtlamayanların bilinen
özellikleri raporlanmalıdır.

2.400 kişilik büyük örneklem rastgele dalgalanmayı azaltabilir; ancak seçim
mekanizmasının oluşturduğu sistematik yanlılığı ortadan kaldırmaz. Bu nedenle
``büyük örneklem temsil edicidir'' sonucu geçerli değildir.

\subsection*{Problem 5: Etik, veri yapısı ve analiz planı}

Bir araştırma ekibi öğrencilerin psikolojik iyi oluş puanlarıyla devamsızlık
bilgilerini birleştirmek istemektedir. Dosyada öğrenci numarası, ad, program,
iyi oluş puanı ve devamsızlık günü vardır. Araştırma başlamadan önce hangi
istatistiksel ve etik kararlar verilmelidir?

\textbf{Çözüm.} İlk olarak araştırma sorusu ve gerçekten gerekli değişkenler
belirlenmelidir. Ad bilgisi analiz için gerekmiyorsa çalışma dosyasından
çıkarılmalıdır. Öğrenci numarası doğrudan kimlik belirlemeye yol açabileceği
için rastgele bir araştırma koduyla değiştirilip eşleştirme anahtarı ayrı ve
güvenli biçimde saklanmalıdır. Erişim yetkisi, saklama süresi, katılımcı onamı
ve gerekli etik kurul/kurum izinleri veri toplanmadan önce planlanmalıdır.

Her satırın tek bir öğrenciyi temsil edip etmediği, aynı öğrencinin birden
fazla kaydının bulunup bulunmadığı, eksik puanların hangi kodla gösterildiği
ve olanaksız değerlerin varlığı kontrol edilmelidir. Örneğin iyi oluş ölçeği
0--100 aralığındaysa 140 değeri ölçüm veya veri giriş hatası olarak
incelenmelidir; doğrudan silinmemelidir.

İlk analizde değişkenlerin dağılımları ve eksik veri örüntüleri özetlenebilir,
ardından iyi oluş ile devamsızlık arasındaki ilişki grafiksel olarak
incelenebilir. Analiz planında hangi gözlemlerin dışlanacağı ve hangi
yöntemlerin kullanılacağı sonuçlar görülmeden önce belirtilmelidir.

\textbf{Raporlama örneği.} ``Kimlik bilgileri analiz dosyasından ayrılmış,
doğrudan tanımlayıcılar araştırma kodlarıyla değiştirilmiştir. Veri temizleme
kuralları analiz öncesinde belirlenmiş; eksik, yinelenen ve ölçüm aralığı
dışındaki değerler kayıt altına alınarak incelenmiştir.''

\section{Literatür verisiyle yazılım uygulamaları}
\label{sec:spss-b01}

\textbf{Amaç ve veri.} İki okulun matematik verisinde bir satırın ve bir
sütunun anlamını belirleyelim. \texttt{b01.csv}, 395 öğrencinin seçilmiş
10 değişkenini içerir. Kaynak ve veri üretim süreci için
\citet{cortez2008data} ve sayfa~\pageref{sec:spss-guide}'teki rehber kullanılır.


\textbf{Ortak veri.} Doğrulanan B01 pilotunda üç yazılım da
\nolinkurl{companion/spss/csv/b01.csv} dosyasındaki aynı kayıtları kullanır.
\texttt{b01.xlsx} aynı verinin alternatif kopyasıdır; bu karşılaştırma
Excel içe aktarmanın doğrulandığı anlamına gelmez. Kodları
\texttt{companion} klasörünü içeren paket kökünden çalıştırın; veri sözlüğü
ve hazırlık için bkz. sayfa~\pageref{sec:spss-guide}.

\subsection{IBM SPSS uygulaması}
\textbf{Başlamadan önce.} B01 CSV pilot paketindeki
\texttt{open-b01.sps} dosyasını çalıştırın; CSV verisi etiketleriyle açılır.
Menü yolu ve aşağıdaki tam komut dosyası alternatif uygulama yollarıdır.

\begin{spssadimlar}
\spssadim{\emph{Variable View} sekmesinde etiket ve ölçme düzeylerini kontrol edin: \texttt{school} ve \texttt{sex} nominal, \texttt{studytime} sıralı, notlar Scale olmalıdır.}
\spssadim{\emph{Analyze $\to$ Descriptive Statistics $\to$ Frequencies} yolunu açın; \texttt{school} ve \texttt{sex} değişkenlerini \emph{Variable(s)} alanına taşıyın.}
\spssadim{\emph{Display frequency tables} seçili kalsın; \emph{OK} ile tabloları alın. Her tablonun geçerli gözlem toplamını kontrol edin.}
\spssadim{\emph{Analyze $\to$ Descriptive Statistics $\to$ Descriptives} yolunda \texttt{age}, \texttt{g1}, \texttt{g2}, \texttt{g3} değişkenlerini seçin.}
\spssadim{\emph{Options} içinde \emph{Mean}, \emph{Std. deviation}, \emph{Minimum}, \emph{Maximum} seçin; \emph{Continue $\to$ OK} ile özetleri alın.}
\end{spssadimlar}

{\raggedright
\noindent\textbf{Komutların görevi.} \texttt{DISPLAY DICTIONARY} veri tanımlarını, \texttt{FREQUENCIES} kategori sayımlarını, \texttt{DESCRIPTIVES} nicel özetleri verir.\par}

\par\noindent\begin{minipage}{\linewidth}
\textbf{Uygulama kodu:} \texttt{b01}. Veri: \texttt{csv/b01.csv};
komut: \texttt{syntax/b01.sps} (\texttt{companion/spss} altında).

% Derleme harici .sps dosyasına bağlı değildir. Eşlikçi kopya:
% companion/spss/syntax/b01.sps
\begin{lstlisting}[style=prismSPSS]
INSERT FILE='companion/spss/syntax/open-b01.sps'
 ERROR=STOP.
DISPLAY DICTIONARY.
FREQUENCIES VARIABLES=school sex.
DESCRIPTIVES VARIABLES=age g1 g2 g3
 /STATISTICS=MEAN STDDEV MIN MAX.
\end{lstlisting}

\end{minipage}\par

\begin{outputguidebox}
\textbf{Çıktıyı okuma ve kontrol.}
\emph{Dictionary} çıktısında \texttt{school} ve \texttt{sex} nominal,
\texttt{studytime} sıralı, notlar ise Scale olarak görünmelidir.
Frekans tablosunda GP için 349, MS için 46; F için 208, M için 187 kayıt
beklenir. Her iki tablonun toplamı 395'tir. \emph{Descriptives} tablosunda
yıl sonu notunun ortalaması \SPMatMean, standart sapması \SPMatSD\ olur.
Bu özetler yalnız gözlenen öğrencileri betimler; 395 satır bulunması
öğrencilerin rastgele seçildiğini veya bağımsız olduğunu kanıtlamaz.
\end{outputguidebox}


\par\noindent\begin{minipage}{\linewidth}
\subsection{R uygulaması}
\label{sec:r-gercek-b01}

\texttt{str} veri türünü, \texttt{is.na} eksikleri, \texttt{table} frekansları gösterir. Nominal ve sıralı değişkenler açıkça tanımlanır.

\begin{lstlisting}[style=prismR]
d <- read.csv("companion/spss/csv/b01.csv")
stopifnot(!anyNA(d))
d$school <- factor(d$school, 1:2, c("GP", "MS"))
d$sex <- factor(d$sex, 1:2, c("F", "M"))
d$studytime <- ordered(d$studytime, levels=1:4)
str(d)
print(colSums(is.na(d)))
print(table(d$school)); print(table(d$sex))
x <- d[c("age", "g1", "g2", "g3")]
print(sapply(x, function(v) c(n=length(v),
  mean=mean(v), sd=sd(v), min=min(v), max=max(v))))
\end{lstlisting}

\textbf{Çıktıyı okuma.} R çıktısında \texttt{n}, \texttt{mean},
\texttt{sd}, \texttt{min} ve \texttt{max} satırlarını aşağıdaki ortak
tablonun sütunlarıyla eşleştirin. Paylaşılan R çıktısındaki frekanslar
ve bu beş özet doğrulanmıştır; eksik değer sayımları paylaşılan dosyada
yer almamaktadır.
\end{minipage}\par

\par\noindent\begin{minipage}{\linewidth}
\subsection{Python uygulaması}
\label{sec:python-b01}

\texttt{info} sütun yapısını, \texttt{isna} eksikleri, \texttt{agg} seçilen özetleri verir. Kategori kodlarının ortalaması alınmaz.

\begin{lstlisting}[style=prismPythonApp]
import pandas as pd
import numpy as np
from scipy import stats

d = pd.read_csv("companion/spss/csv/b01.csv")
assert not d.isna().any().any()
d["school"] = pd.Categorical(d.school.map({1:"GP", 2:"MS"}))
d["sex"] = pd.Categorical(d.sex.map({1:"F", 2:"M"}))
d["studytime"] = pd.Categorical(
    d.studytime, categories=[1, 2, 3, 4], ordered=True)
d.info()
print(d.isna().sum())
print(d.school.value_counts()); print(d.sex.value_counts())
print(d[["age", "g1", "g2", "g3"]].agg(
    ["count", "mean", "std", "min", "max"]))
\end{lstlisting}

\textbf{Çıktıyı okuma.} Python çıktısı 395 satır ve 10 sütun gösterir;
her sütunda eksik değer sayısı sıfırdır. \texttt{count}, \texttt{mean},
\texttt{std}, \texttt{min} ve \texttt{max} satırlarını ortak tabloyla
karşılaştırın. Bu uygulama betimseldir; bir hipotez testi veya $p$ değeri
üretmez.
\end{minipage}\par

\subsection{Sonuçların karşılaştırılması ve ortak rapor}
12 Eylül 2026'da paylaşılan SPSS ekran görüntüleri ile R ve Python
metin çıktıları karşılaştırılmıştır. R ve Python'daki 20 betimsel hücre
gösterilen altı ondalık basamakta eşleşmiş; SPSS'in yuvarlanmış
değerleriyle de uyum sağlanmıştır. Aşağıdaki tablolar bu çıktıların
ortak, kitap düzenine uyarlanmış özetidir; ekran görüntüsü değildir.

\begin{table}[htbp]
\centering
\caption{B01 verisinde okul ve kaynak cinsiyet kodlarının dağılımı}
\label{tab:b01-dogrulanmis-frekans}
\begin{tabular}{@{}llrr@{}}
\toprule
Değişken & Kategori & Frekans & Yüzde \\
\midrule
Okul & GP & 349 & 88,4 \\
     & MS & 46 & 11,6 \\
\midrule
Cinsiyet kodu & F & 208 & 52,7 \\
              & M & 187 & 47,3 \\
\bottomrule
\end{tabular}
\par\smallskip
{\footnotesize Her değişken için toplam 395 kayıt ve yüzde 100'dür.
Yüzdeler kategori sayısının 395'e bölünmesiyle hesaplanmıştır.\par}
\end{table}

\begin{table}[htbp]
\centering
\caption{B01 verisinin doğrulanmış betimsel istatistikleri}
\label{tab:b01-dogrulanmis-betimsel}
\small
\begin{tabular}{@{}lrrrrr@{}}
\toprule
Değişken & $n$ & Ortalama & Std. sapma & Min. & Maks. \\
\midrule
Yaş & 395 & 16,70 & 1,276 & 15 & 22 \\
\texttt{g1} & 395 & 10,91 & 3,319 & 3 & 19 \\
\texttt{g2} & 395 & 10,71 & 3,762 & 0 & 19 \\
\texttt{g3} & 395 & 10,42 & 4,581 & 0 & 20 \\
\bottomrule
\end{tabular}
\par\smallskip
{\footnotesize Ortalamalar iki, örneklem standart sapmaları üç ondalık
basamağa yuvarlanmıştır. SPSS'te ortak geçerli gözlem sayısı
(Valid N, listwise) 395'tir.\par}
\end{table}

\begin{outputguidebox}
Tablo~\ref{tab:b01-dogrulanmis-frekans}, öğrencilerin okullara eşit
dağılmadığını gösterir. Tablo~\ref{tab:b01-dogrulanmis-betimsel} içinde yıl
sonu notu ortalaması 10,42, standart sapması 4,581 puandır. Standart sapma
öğrenci notlarının yayılımını, ortalama ise merkezini özetler; bu iki
sayı aynı bilgiyi vermez. Üç yazılımın uyumu hesapların tutarlılığını
destekler; örneklemin evreni temsil ettiğini veya nedenselliği kanıtlamaz.
Dosya türü ile ölçme düzeyi farklı kavramlardır.
\end{outputguidebox}

\textbf{Raporlama örneği (doğrulanmış çıktılarla).}
``İki okuldan 395 öğrenci kaydının 349'u (\%88,4) GP, 46'sı (\%11,6)
MS okuluna aittir. Öğrencilerin yaş ortalaması 16,70 yıl
($s=1{,}276$), yıl sonu matematik notu ortalaması 10,42 puandır
($s=4{,}581$). Yıl sonu notları 0 ile 20 arasında değişmektedir.
Bu betimleme daha geniş bir öğrenci evrenine doğrudan genellenmemiştir.''


\textbf{Görev.} Raporlama örneğine hedef evren tanımını ekleyin; okul kodlarının ortalamasının neden raporlanmayacağını açıklayın.

\section{R ile çözümlü uygulama}
\label{sec:r-b01}

\textbf{Soru.} Bölümdeki beş öğrencilik öğretim örneğini R'de oluşturun.
Gözlem birimini, değişken türlerini, program frekanslarını ve ortalama başarıyı
belirleyin. Bu küçük veri gerçek bir araştırmadan alınmış değildir.

\begin{lstlisting}[style=prismR]
veri <- data.frame(
  devam_saati = c(8, 10, 7, 12, 9),
  basari = c(68, 75, 64, 83, 72),
  program = factor(c("A", "A", "B", "B", "A"))
)
str(veri)
summary(veri)
table(veri$program)
prop.table(table(veri$program))
colSums(is.na(veri))
c(ogrenci = nrow(veri),
  ortalama_devam = mean(veri$devam_saati),
  ortalama_basari = mean(veri$basari))
\end{lstlisting}

\begin{outputguidebox}
\textbf{Çözüm ve kontrol değerleri.} Her satır bir öğrencidir; veri
5 satır ve 3 değişken içerir. \texttt{devam\_saati} ve \texttt{basari}
nicel, \texttt{program} nominal kategoriktir. A programında 3, B programında
2 öğrenci vardır. Ortalama devam 9,2 saat, ortalama başarı 72,4 puandır;
hiçbir sütunda eksik değer yoktur.
\end{outputguidebox}

\textbf{Yorum.} \texttt{factor} program kodlarının kategori olarak
saklanmasını sağlar. Bir program etiketinin ortalamasını almak anlamlı
değildir. Beş satırın özetlenmesi, devamın başarıya neden olduğunu veya
öğrencilerin bir evreni temsil ettiğini göstermez.

\textbf{Pekiştirme ve yanıt.} Programların örneklemdeki paylarını
\texttt{prop.table} çıktısından okuyun: A için 0,60, B için 0,40 elde edilir. Bunlar evren oranı değil örneklem oranlarıdır.

\section{Bölüm özeti}

\begin{itemize}
  \item İstatistiksel çalışma araştırma sorusu ve veri üretim tasarımıyla başlar.
  \item Değişken türü, kullanılabilecek grafik ve yöntemleri sınırlar.
  \item Örneklemden evrene yapılan her genelleme belirsizlik ve olası yanlılık içerir.
  \item Etik, şeffaflık ve nedensellik sınırları istatistiksel raporun parçasıdır.
\end{itemize}

\section*{Alıştırmalar}

\begin{enumerate}
  \item Bir eğitim araştırması sorusu yazın; hedef evreni, gözlem birimini,
  açıklayıcı değişkeni ve yanıt değişkenini belirtin.
  \item Okul türü, yaş, öğrenci numarası, başarı sırası ve sınav süresinin
  ölçme düzeylerini sınıflandırın; her biri için uygun bir özet önerin.
  \item Bir ilçedeki öğrencilerin ortalama okuma süresini inceleyen çalışmada
  parametre ve istatistiğe birer örnek verin.
  \item Bir çevrim içi ankette yalnız gönüllü katılımcıların kullanılmasının
  doğurabileceği iki yanlılık kaynağını açıklayın.
  \item Tabakalı örneklemenin basit rastgele örneklemeye tercih edilebileceği
  bir araştırma durumu tasarlayın.
  \item Korelasyon gözlenen bir çalışmada nedensel sonuç çıkarılmasını
  engelleyebilecek en az iki karıştırıcı değişken önerin.
  \item Her satırın bir sınav girişini gösterdiği veri dosyasında aynı
  öğrencinin iki sınava girmesinin analiz birimi bakımından doğuracağı sorunu
  açıklayın.
  \item Eksik değerlerin 0 koduyla kaydedilmesinin ortalama ve dağılım üzerinde
  nasıl yanıltıcı bir etki oluşturabileceğini tartışın.
  \item Büyük bir örneklemin neden kötü bir örnekleme çerçevesini telafi
  edemeyeceğini bir örnekle açıklayın.
  \item Bir öğrenci veri seti için veri toplama, temizleme, analiz ve etik
  raporlama adımlarını içeren kısa bir araştırma planı hazırlayın.
\end{enumerate}